% ─────────────────────────────────────────────────────────────────
%  Draft of §IV.B (relaxation oscillator) — to replace existing
%  §IV.C (Mechanism of the LCST-front oscillation, lines 487–555 of
%  main.tex).  Uses figures B2 (slow_manifold), B3 (period_scaling),
%  and the existing time-series / heatmap panels of current Fig.4.
% ─────────────────────────────────────────────────────────────────

\subsection{A relaxation oscillator on a spatially-generated slow
manifold}
\label{sec:lcst_front_mechanism}

The LCST-front cycle is a relaxation oscillator with two distinguishing
features: (i)~the fast variable is mechanical ($J$), the slow
variables are thermal and chemical ($\theta$, $u$); and (ii)~the
slow manifold is generated by the spatial structure of the gel rather
than by an autonomous 0D nonlinearity.  Each cycle traverses two
stable branches of this manifold connected by fast jumps, and the
period is dominated by the slowest of the four phases.

\paragraph{Fast--slow decomposition.}  At the working point the
mechanical relaxation timescale set by the dimensionless mobility
$\Ms_0$ is roughly two orders of magnitude shorter than the thermal
($1/\Bi_T$) and reactant-supply ($1/\Bi_c$) timescales.  Numerically,
the speed $|d(\theta_\mathrm{surf}, J_\mathrm{surf})/dt|$ along the
limit cycle [Fig.~\ref{fig:fig2_mechanism}(c)] varies by a factor of
$\sim\!600$ between the slow drifts and the fast jumps.  We may
therefore treat $J$ as a fast variable and follow the standard
relaxation-oscillator construction.

\paragraph{Spatially generated slow manifold.}
Quasi-static balance of the surface mechanical equation with the
Robin boundary condition $\Bi_\mu(\mu - \mu_b) = 0$ gives the locus
\begin{equation}
  \mu(J,\theta) \;=\; \mu_b
  \label{eq:slow_manifold}
\end{equation}
in the $(\theta, J)$ plane.  The Flory--Huggins--Rehner free energy
is non-monotonic in $J$ for sufficiently large $\chi_\mathrm{eff}=
\chi_\infty + S_\chi\theta + \chi_1 \phi$ (above the LCST), so this
locus has the canonical S-shape
[Fig.~\ref{fig:fig2_mechanism}(c)]: an upper \emph{swollen} branch
($J \gtrsim 0.8$) for $\theta \le \theta_\mathrm{up}$, a lower
\emph{collapsed} branch ($J \lesssim 0.2$) for
$\theta \ge \theta_\mathrm{lo}$, and a saddle (unstable middle
branch) connecting the two folds. At the working point,
$\theta_\mathrm{lo}\!\simeq\!0.85$ and
$\theta_\mathrm{up}\!\simeq\!3.18$, defining a bistable region
$\theta\in[\theta_\mathrm{lo},\theta_\mathrm{up}]$ in surface
temperature.  Crucially, this slow manifold lives at the surface and
is enforced by the surface BC; it has no analog in a 0D
homogeneous-state reduction, where $\mu = \mu_b$ would have to be
imposed everywhere simultaneously together with the heat and reactant
balances (Sec.~\ref{sec:lsa_vs_nl}).

\paragraph{Cycle traversal.}  The PDE limit cycle, projected onto
$(\theta_\mathrm{surf}, J_\mathrm{surf})$, traces the slow manifold's
two stable branches connected by near-vertical fast jumps
[Fig.~\ref{fig:fig2_mechanism}(c)].  We label four canonical phases:

\begin{enumerate}[leftmargin=*,label=\arabic*.]
\item \textbf{Ignition (slow, on swollen branch).}  The surface starts
  at the cold end of the swollen branch ($\theta\!\simeq\!1.5$,
  $J\!\simeq\!1.4$) with abundant reactant
  ($u_\mathrm{surf}\!\simeq\!1$).  A small thermal fluctuation
  amplifies the Arrhenius factor; the heat source $\Da\,J\hat R$
  exceeds the surface heat loss $\Bi_T\theta_\mathrm{surf}$ and
  $\theta_\mathrm{surf}$ rises monotonically.
\item \textbf{Collapse (fast).}  As $\theta_\mathrm{surf}$ exceeds the
  upper-fold value, the swollen branch ceases to exist; $J$ falls
  rapidly along the fast manifold from $J\!\sim\!0.8$ to
  $J\!\sim\!0.15$ over a small fraction of the cycle period.
\item \textbf{Quenching (slow, on collapsed branch).}  With
  $\varphi\!\sim\!1$ in the collapsed shell the diffusion factor
  $D(\varphi) \propto (1-\varphi)^2$ drops by two orders of magnitude,
  throttling reactant supply from the bath.  Boundary-layer reactant
  is consumed by the residual reaction; with the heat source gone,
  $\theta_\mathrm{surf}$ relaxes towards the bath temperature on the
  surface heat-loss timescale $1/\Bi_T$.  This is generally the
  longest phase of the cycle.
\item \textbf{Re-swelling (fast).}  When $\theta_\mathrm{surf}$ falls
  below the lower-fold value, the collapsed branch terminates and $J$
  jumps back up to the swollen branch.  The diffusion barrier
  reopens, the bath rapidly resupplies reactant, and the cycle
  returns to phase~1.
\end{enumerate}

The slow drifts (phases 1 and 3) and the fast jumps (phases 2 and 4)
are clearly resolved both as time-series plateaus and ramps
[Fig.~\ref{fig:fig2_mechanism}(b)] and as differently-shaded
segments of the limit cycle in the slow-manifold plane
[Fig.~\ref{fig:fig2_mechanism}(c)].

\paragraph{Self-limiting collapse.}  The same $\varphi$-dependent
diffusion barrier that quenches the reaction at the surface during
phase~3 also prevents collapse from invading the gel interior.  The
LCST front penetrates only to
$\xi_\mathrm{LCST}\!\simeq\!0.885$ at the working point; the inner
$88.5\%$ of the slab never crosses the LCST and remains swollen
throughout the cycle.  This barrier explains why the classifier-
allowed regimes \emph{globally collapsed steady} and \emph{globally
collapsing oscillation} are absent from $1{,}000$ default-IC
simulations spanning the parameter space (Sec.~\ref{sec:phasediagram}):
once any portion of the gel collapses, the inner region is shielded
from further reactant supply and cannot ignite. The same mechanism
is, in this sense, a self-limiting protective feature of this class
of thermo-LCST gel against runaway global collapse despite the
strongly exothermic Arrhenius
kinetics~\cite{FrankKamenetskii1969,Uppal1974}.

\paragraph{Period structure.}  Across the LCST-front regime
[Fig.~\ref{fig:fig2_mechanism}(e)], the cycle period $T_\mathrm{PDE}$
spans roughly $6$--$15$, with a minimum of $\sim\!8$ near the
``sweet-spot'' $S_\chi\!\simeq\!0.75$ and divergence at the upper-
$S_\chi$ boundary where the cycle disappears through a saddle--node-
on-invariant-circle (SNIC)-like bifurcation
(Sec.~\ref{sec:lsa_vs_nl}).  The cooling timescale $1/\Bi_T$ sets the
order of magnitude of the period (prefactor $\mathcal{O}(1)$
absorbing the ignition delay and the spatial barrier-release time),
but the period is not a simple power law of $\Bi_T$ alone because
phase~3 also depends on the residual interior heat conduction and on
the timescale at which the boundary BC restores
$u_\mathrm{surf}$ once the barrier reopens.
